\documentclass[12pt]{article}

\usepackage[spanish]{babel}
\usepackage[utf8]{inputenc}
\usepackage[T1]{fontenc}
\usepackage{geometry}
\usepackage{graphicx}
\usepackage{hyperref}
\usepackage{array}
\usepackage{xcolor}
\usepackage{longtable}
\usepackage{enumitem}
\usepackage{titlesec}

\geometry{
    letterpaper,
    margin=2.5cm
}

\hypersetup{
    colorlinks=true,
    linkcolor=blue,
    urlcolor=blue
}

\title{\textbf{Documentación del Frontend del Sistema de Reservas del Gimnasio UNAL}}
\author{Equipo No Pain No Main}
\date{\today}

\begin{document}

\maketitle

\section{Introducción}

Este documento presenta el trabajo realizado en el frontend del sistema de reservas de máquinas del gimnasio de la Universidad Nacional de Colombia, sede Bogotá. El desarrollo se llevó a cabo con Vue.js y Vite, utilizando una arquitectura modular orientada a separar responsabilidades, mejorar la mantenibilidad y preparar el proyecto para una futura integración con el backend.

El frontend fue reconstruido en una rama independiente llamada:\\
\texttt{feature/frontend-architecture}, con el objetivo de trabajar de forma ordenada sin afectar la rama principal del repositorio.

Dado que el backend y la base de datos ya existían, pero aún no contaban con endpoints completamente implementados, se decidió construir la interfaz con datos simulados o \textit{mocks}. De esta manera, las pantallas pudieron probarse visual y funcionalmente antes de conectar la aplicación con servicios reales.

\section{Objetivo del trabajo}

El objetivo principal fue construir una interfaz frontend funcional, clara y preparada para integrarse con el backend del sistema de reservas.

Los objetivos específicos fueron:

\begin{itemize}
    \item Reconstruir el proyecto frontend desde cero usando Vue.js.
    \item Definir una arquitectura por capas que facilitara el mantenimiento del código.
    \item Implementar las pantallas principales del sistema definidas en los casos de uso y mockups.
    \item Utilizar datos simulados mientras se completaban los endpoints del backend.
    \item Preparar la aplicación para una futura conexión con la API.
    \item Documentar la estructura y las decisiones técnicas tomadas durante el desarrollo.
\end{itemize}

\section{Tecnologías utilizadas}

Para el desarrollo del frontend se utilizaron las siguientes tecnologías:

\begin{itemize}
    \item \textbf{Vue.js 3}: framework principal para construir la interfaz de usuario.
    \item \textbf{Vite}: herramienta de desarrollo y compilación del proyecto.
    \item \textbf{Vue Router}: gestión de rutas y navegación entre pantallas.
    \item \textbf{Pinia}: librería preparada para el manejo del estado global.
    \item \textbf{Axios}: cliente HTTP para futuras peticiones al backend.
    \item \textbf{Git y GitHub}: control de versiones y trabajo colaborativo por ramas.
    \item \textbf{PowerShell}: terminal utilizada para ejecutar comandos y gestionar el proyecto.
\end{itemize}

\section{Creación de la rama de trabajo}

El proceso de desarrollo comenzó con la clonación del repositorio:

\begin{verbatim}
git clone https://github.com/No-pain-No-main/project-frontend.git
\end{verbatim}

Posteriormente se ingresó a la carpeta del proyecto:

\begin{verbatim}
cd project-frontend
\end{verbatim}

Y se creó una rama independiente para trabajar sin afectar la línea principal del proyecto:

\begin{verbatim}
git checkout -b feature/frontend-architecture
\end{verbatim}

\section{Reconstrucción del proyecto Vue}

Al revisar el repositorio se encontró que varios archivos base del proyecto Vue habían sido eliminados en commits anteriores, entre ellos:

\begin{itemize}
    \item \texttt{package.json}
    \item \texttt{vite.config.js}
    \item \texttt{index.html}
    \item carpeta \texttt{src/}
    \item carpeta \texttt{public/}
\end{itemize}

Por esta razón, se decidió reconstruir la base del proyecto desde cero con Vite mediante los siguientes comandos:

\begin{verbatim}
npm create vite@latest . -- --template vue
npm install
npm install vue-router@4 axios pinia
\end{verbatim}

Una vez creado el proyecto, se verificó su funcionamiento ejecutando:

\begin{verbatim}
npm run dev
\end{verbatim}

El servidor local quedó disponible en la ruta:

\begin{verbatim}
http://localhost:5173/
\end{verbatim}

\section{Arquitectura por capas}

Para organizar el frontend, se adoptó una arquitectura por capas que permitiera separar responsabilidades y facilitar el mantenimiento del proyecto. Esta estructura ayuda a que cada parte del sistema tenga un propósito claro y sea más fácil de extender.

La estructura principal del proyecto quedó organizada de la siguiente manera:

\begin{verbatim}
src/
|-- components/
|-- layouts/
|-- mocks/
|-- router/
|-- services/
|-- stores/
|-- styles/
|-- views/
|-- App.vue
|-- main.js
\end{verbatim}

\subsection{Capa de vistas}

La carpeta \texttt{views} contiene las pantallas completas del sistema. Cada vista representa una página o módulo principal del proyecto.

Se organizaron en categorías como:

\begin{itemize}
    \item pantallas de autenticación,
    \item pantallas del estudiante,
    \item pantallas del administrador.
\end{itemize}

\subsection{Capa de componentes}

La carpeta \texttt{components} contiene elementos reutilizables de la interfaz, como formularios, tarjetas, botones y otros bloques visuales que pueden usarse en varias vistas.

\subsection{Capa de layouts}

Los layouts permiten definir estructuras visuales comunes para varias páginas. Se crearon layouts específicos para la autenticación, el estudiante y el administrador.

\subsection{Capa de rutas}

La carpeta \texttt{router} centraliza la configuración de navegación. Aquí se definieron rutas públicas y rutas protegidas, además de la lógica para dirigir a cada usuario según su rol.

\subsection{Capa de mocks}

Debido a la falta de endpoints del backend, se implementó una capa de datos simulados en la carpeta \texttt{mocks}. Esto permitió probar el comportamiento general del sistema sin depender todavía de una API real.

\subsection{Capa de servicios}

La carpeta \texttt{services} quedó preparada para encapsular la comunicación con el backend en una etapa posterior. Esta decisión permite que los componentes y las vistas no dependan directamente de la lógica de conexión.

\subsection{Capa de estado}

La carpeta \texttt{stores} se destinó al manejo del estado global de la aplicación mediante Pinia, lo que deja la puerta abierta para manejar autenticación, usuario activo, reservas y otras variables compartidas.

\section{Pantallas implementadas}

Se desarrollaron varias pantallas principales del sistema, orientadas a cubrir los flujos más importantes del negocio.

\subsection{Inicio de sesión}

Se implementó la pantalla de autenticación con campos para ingresar el usuario y la contraseña. Actualmente funciona de forma simulada, pero su estructura está lista para conectarse con un servicio de autenticación real.

\subsection{Registro de usuario}

Se desarrolló la vista de registro con campos básicos para crear una cuenta de estudiante, incluyendo datos personales y credenciales de acceso.

\subsection{Panel del estudiante}

Se creó un panel para el estudiante con acceso a información principal y a las opciones más relevantes de uso del sistema, como consultar reservas, ver máquinas disponibles y gestionar su perfil.

\subsection{Reservar máquina}

Se implementó una vista para reservar máquinas, mostrando tarjetas de equipos con información relacionada a su estado y permitiendo seleccionar una franja horaria para confirmar una reserva.

\subsection{Mis reservas}

Se creó una vista donde el estudiante puede visualizar sus reservas activas y cancelarlas cuando sea necesario.

\subsection{Panel del administrador}

Se diseñó un panel administrativo con opciones para gestionar máquinas, usuarios, reservas y reportes, con el fin de facilitar la administración del sistema.

\subsection{Gestión de máquinas}

Se implementaron pantallas para registrar máquinas y modificar su estado de disponibilidad, lo que permite simular el control de los equipos del gimnasio.

\subsection{Historial de reservas}

Se creó una vista para consultar reservas de forma general, con filtros y tablas que facilitan la revisión de la información.

\subsection{Reportes administrativos}

Se implementó una vista para mostrar reportes base y estadísticas del uso del sistema, con una interfaz preparada para mejorar en una segunda fase.

\section{Relación con los casos de uso}

Las pantallas desarrolladas corresponden directamente a los principales casos de uso del sistema:

\begin{longtable}{|p{4cm}|p{4cm}|p{6cm}|}
\hline
\textbf{Caso de uso} & \textbf{Pantalla implementada} & \textbf{Descripción} \\
\hline
Registrar usuario & Registro de usuario & Permite crear una cuenta para el estudiante. \\
\hline
Iniciar sesión & Login & Permite entrar al sistema según el rol del usuario. \\
\hline
Reservar máquina & Reservar máquina & Permite seleccionar una máquina y una franja horaria. \\
\hline
Cancelar reserva & Mis reservas & Permite cancelar reservas activas. \\
\hline
Registrar máquina & Gestión de máquinas & Permite agregar nuevas máquinas al sistema. \\
\hline
Modificar estado de máquina & Gestión de máquinas & Permite cambiar el estado de disponibilidad de un equipo. \\
\hline
Consultar historial & Historial de reservas & Permite revisar reservas pasadas y presentes. \\
\hline
Generar reportes & Reportes administrativos & Permite visualizar estadísticas de uso. \\
\hline
\end{longtable}

\section{Decisiones importantes}

\subsection{Uso de datos simulados}

Se decidió trabajar con datos simulados porque el backend aún no contaba con los endpoints necesarios para integrar el sistema de forma completa. Esta decisión permitió avanzar en el frontend sin detener el desarrollo.

\subsection{Separación de responsabilidades}

Se evitó concentrar toda la lógica en las vistas. En su lugar, se separaron responsabilidades en vistas, componentes, servicios, mocks y stores, lo que facilita el mantenimiento y la escalabilidad del proyecto.

\subsection{Trabajo en rama independiente}

El desarrollo se realizó en una rama separada para evitar afectar la rama principal del repositorio y permitir una revisión más controlada de los cambios.

\section{Comandos principales utilizados}

\begin{verbatim}
git clone https://github.com/No-pain-No-main/project-frontend.git
cd project-frontend
git checkout -b feature/frontend-architecture
npm create vite@latest . -- --template vue
npm install
npm install vue-router@4 axios pinia
npm run dev
git add .
git commit -m "mensaje del commit"
\end{verbatim}

\section{Estado actual del proyecto}

Actualmente el frontend cuenta con una estructura sólida, navegación funcional, vistas principales implementadas y un flujo de uso que permite probar el sistema con datos simulados. El proyecto ya está preparado para avanzar hacia una integración más real con el backend.

\section{Trabajo pendiente}

Aún quedan tareas importantes por completar, entre ellas:

\begin{itemize}
    \item implementar servicios reales para consumir la API del backend,
    \item conectar login y registro con la base de datos,
    \item persistir reservas y cambios de estado,
    \item generar reportes a partir de datos reales,
    \item mejorar la validación de formularios,
    \item ampliar la cobertura de pruebas y revisión del código.
\end{itemize}

\section{Conclusión}

Se reconstruyó el frontend del sistema de reservas usando Vue.js y una arquitectura por capas clara, organizada y escalable. El proyecto quedó con una base sólida para la interacción del usuario, con vistas adecuadas para estudiantes y administradores, además de una estructura lista para crecer hacia una integración completa con el backend.

Este avance representa una base sólida para continuar desarrollando el sistema de manera más completa y profesional.

\end{document}
